\chapter{Introduction}\label{ch:m-introduction}
\section{Running example and research gap}
A cyber-physical monitoring system receives video, audio, and vibration streams. During a critical event, the camera is partially occluded and the audio stream is delayed. A conventional fusion model may still produce a confident prediction even though two inputs are unreliable. The research problem is therefore not only to improve nominal accuracy, but to preserve useful behaviour and calibrated uncertainty when modality quality changes.

Multimodal learning provides mechanisms for combining heterogeneous information \cite{goodfellow2016deep}, while trustworthy deployment requires explicit treatment of reliability, transparency, and risk \cite{nist2023airmf}. This thesis focuses on the gap between high performance under complete data and dependable behaviour under missing or degraded modalities.

\begin{researchquestion}[RQ1 -- Robustness]
How much predictive performance can a reliability-aware fusion model preserve when one sensor stream is noisy, delayed, or unavailable?
\end{researchquestion}

\begin{researchquestion}[RQ2 -- Trust]
Which design and evaluation choices make the model's confidence, failure conditions, and computational cost transparent to operators and auditors?
\end{researchquestion}

\section{Central thesis claim and contributions}
\begin{resultbox}[Central thesis claim]
Estimating modality reliability before fusion can improve robustness and calibration under controlled sensor degradation without exceeding the latency budget of the target platform.
\end{resultbox}

The thesis makes four refutable contributions:
\begin{enumerate}
  \item a problem formulation with explicit assumptions, degradation conditions, and success criteria in Chapter~\ref{ch:m-problem};
  \item a modular reliability-aware architecture and objective in Chapter~\ref{ch:m-methodology};
  \item a reproducible evaluation of accuracy, calibration, latency, robustness, uncertainty, and failure cases in Chapter~\ref{ch:m-results};
  \item a synthesis that compares assumptions with prior work and checks every central claim against evidence in Chapters~\ref{ch:m-related} and~\ref{ch:m-discussion}.
\end{enumerate}

\chapter{Problem Formulation and Running Example}\label{ch:m-problem}
\section{Inputs, degradation, and desired behaviour}
Let the system receive modalities $x_1,\ldots,x_M$. Each modality may be available, missing, delayed, or corrupted according to a predefined test condition. The model produces a prediction and an uncertainty estimate. The desired behaviour is graceful degradation: confidence should fall when evidence quality falls, and the system should avoid a large, unexplained performance collapse.

\section{Success criteria and scope}
The example evaluates four properties: macro F1, expected calibration error, robustness under missing modalities, and median end-to-end latency. It assumes synchronised labels, a fixed set of sensors, and degradation conditions generated by a documented protocol. It does not claim robustness to every adversarial manipulation or to unseen sensor types.

\begin{table}[htbp]
\centering
\caption{Illustrative criteria for the master thesis example.}
\label{tab:m-criteria}
\begin{tabularx}{\textwidth}{YYY}
\toprule
Property & Example criterion & Evidence\\
\midrule
Nominal quality & Macro F1 above the strongest baseline & Repeated held-out evaluation\\
Missing-modality robustness & Relative loss below 15\% at 30\% missing inputs & Controlled modality dropout\\
Calibration & Expected calibration error below 0.05 & Reliability diagram and scalar metric\\
Latency & Median below 35~ms & End-to-end timing on target hardware\\
\bottomrule
\end{tabularx}
\end{table}

\chapter{Proposed Method}\label{ch:m-methodology}
\section{Whiteboard intuition}
Each sensor encoder proposes a representation and reports a reliability score. The fusion module gives less influence to a modality whose signal is missing or inconsistent with the others. The prediction head returns both a decision and uncertainty information. This intuition comes before equations so the reader knows what the formal terms are meant to achieve.

\begin{figure}[htbp]
\centering
\pipelinebox{Data quality\\validation}\pipelinearrow
\pipelinebox{Modality\\encoders}\pipelinearrow
\pipelinebox{Reliability-aware\\fusion}\pipelinearrow
\pipelinebox{Prediction and\\uncertainty}
\caption{Illustrative reliability-aware multimodal architecture.}
\label{fig:m-architecture}
\end{figure}

\section{Formal objective}
For modality representation $z_m=f_m(x_m)$ and reliability estimate $r_m$, the fusion function combines weighted evidence before prediction. A simplified objective combines task loss, calibration regularisation, and consistency under modality dropout:
\begin{equation}
\mathcal{L}=\mathcal{L}_{\mathrm{task}}+\lambda_{c}\mathcal{L}_{\mathrm{cal}}+\lambda_{m}\mathcal{L}_{\mathrm{missing}}.
\end{equation}
Define every term, state how hyperparameters are selected, and separate choices made during exploration from the final confirmatory protocol.

\section{Reproducible implementation}
Record dataset versions, preprocessing, split generation, hyperparameters, random seeds, hardware, software libraries, stopping criteria, and model-selection rules. Store final configurations and a command that rebuilds the principal tables and figures.

\begin{keypoint}[Evidence for the method]
An architecture diagram explains structure but does not prove the design claim. Add ablations that remove reliability weighting, calibration regularisation, and modality dropout so that each claimed mechanism has corresponding evidence.
\end{keypoint}

\chapter{Experimental Evidence}\label{ch:m-results}
\section{Protocol}
Compare against stated baselines using identical data splits and evaluation budgets. Report repeated runs and uncertainty. Predefine missing-modality conditions, including complete loss, random dropout, delay, and noise, rather than selecting only conditions that favour the proposed model.

\section{Headline observations}
\begin{center}
\begin{minipage}{0.30\textwidth}\metric{94.2\%}{macro F1}{ISISBlue}\end{minipage}\hfill
\begin{minipage}{0.30\textwidth}\metric{0.031}{calibration error}{ISISLightBlue}\end{minipage}\hfill
\begin{minipage}{0.30\textwidth}\metric{27 ms}{median latency}{ISISDeepBlue}\end{minipage}
\end{center}

\begin{figure}[htbp]
\centering
\begin{minipage}{0.84\textwidth}
{\sffamily\bfseries\color{ISISDeepBlue}Macro F1 with 30\% missing modalities\par\vspace{3mm}}
\demobar{Baseline}{55mm}{69.0\%}{ISISMuted}
\demobar{Proposed}{68mm}{85.1\%}{ISISBlue}
\end{minipage}
\caption{Illustrative robustness comparison; replace with real estimates and uncertainty.}
\label{fig:m-results}
\end{figure}

\section{Ablation evidence}
\begin{table}[htbp]
\centering
\caption{Illustrative ablation table linking mechanisms to claims.}
\label{tab:m-ablation}
\begin{tabularx}{0.88\textwidth}{Yrrr}
\toprule
Configuration & Macro F1 & F1 at 30\% missing & Calibration error\\
\midrule
Full method & 94.2\% & 85.1\% & 0.031\\
Without reliability weights & 93.8\% & 77.4\% & 0.058\\
Without calibration term & 94.0\% & 84.8\% & 0.091\\
Without modality dropout & 94.3\% & 73.2\% & 0.046\\
\bottomrule
\end{tabularx}
\end{table}

\begin{resultbox}[Main interpretation]
The illustrative ablations suggest that reliability weighting and modality dropout contribute most to robustness, while calibration regularisation mainly affects confidence quality. A real thesis should report uncertainty and verify that these differences are not artefacts of one split or seed.
\end{resultbox}

\section{Failure cases and uncertainty}
Analyse examples where the model is confidently wrong, where reliability scores disagree with actual sensor quality, and where latency exceeds the budget. Report the frequency and conditions of each failure. Negative results define the scope of the central claim and may motivate a safer abstention mechanism.

\chapter{Related Work and Positioning}\label{ch:m-related}
The detailed comparison appears after the method so that technical differences have a clear purpose. Organise the literature by assumptions, degradation model, uncertainty treatment, and deployment constraints instead of presenting a chronological catalogue.

\begin{table}[htbp]
\centering
\caption{Illustrative comparison of multimodal approaches.}
\label{tab:m-related}
\begin{tabularx}{\textwidth}{Ycccc}
\toprule
Approach & Fusion & Missing data & Calibration & Edge-ready\\
\midrule
Early fusion & Input & Limited & No & Partial\\
Late fusion & Decision & Moderate & Optional & Yes\\
Reliability-gated fusion & Hybrid & Explicit & Optional & Partial\\
This thesis example & Hybrid & Explicit & Required & Yes\\
\bottomrule
\end{tabularx}
\end{table}

Describe earlier methods accurately and acknowledge their strengths. State whether the thesis changes the model, the evaluation protocol, the deployment setting, or only the combination of those elements. Avoid claiming novelty from a missing citation search alone.

\chapter{Discussion and Validity}\label{ch:m-discussion}
\section{Claim-evidence audit}
The central robustness claim is supported by controlled modality-loss experiments and ablations. The calibration claim is supported by scalar error and reliability plots. The latency claim is supported only on the specified hardware. The transparency claim requires documentation, traceable configurations, and failure analysis in addition to model metrics.

\section{Threats to validity}
\textbf{Internal validity:} implementation bugs, unequal tuning budgets, or leakage could explain the result. \textbf{Construct validity:} synthetic modality dropout may not represent field failures. \textbf{External validity:} one dataset and platform limit generalisation. \textbf{Conclusion validity:} too few runs or selective reporting can exaggerate differences. State the mitigation used for each threat and the residual risk.

\begin{keypoint}[Separate evidence from inference]
Write measured observations first, then interpretation, then broader implications. Mark hypotheses for future work as hypotheses rather than conclusions established by the current experiments.
\end{keypoint}

\chapter{Conclusions and Outlook}
RQ1 is answered conditionally: under the documented degradation conditions, the illustrative method preserves more macro F1 than the baseline. RQ2 is addressed through calibrated confidence, modular design, reproducible configurations, and explicit failure analysis, but operator trust still requires human-centred evaluation.

\section{Limitations and future directions}
The example does not establish robustness to adversarial manipulation, new sensor types, long-term distribution shift, privacy attacks, or field interaction with operators. Future work should test online adaptation, formal uncertainty guarantees, privacy-aware learning, deployment across devices, and human-centred evaluation of explanations and abstentions.
